\documentclass[12pt, letterpaper]{article}
\usepackage[utf8]{inputenc}
\usepackage[spanish]{babel}
\usepackage{graphicx}
\usepackage{geometry}
\usepackage{longtable}
\usepackage{array}
\usepackage{booktabs}
\usepackage{float}
\usepackage{caption}
\usepackage{hyperref}
\hypersetup{
    colorlinks=true,
    linkcolor=blue,
    urlcolor=blue,
}
\usepackage{url}
\geometry{top=2.5cm, bottom=2.5cm, left=3cm, right=3cm}

\begin{document}

% ============================================
% PORTADA (Ya la tienes, la incluyo por completitud)
% ============================================
\begin{titlepage}
    \centering
    \vspace*{1.5cm}
    
    {\Large \textbf{Universidad del Valle}}\\
    {\Large Facultad de Ingeniería}\\
    {\Large Escuela de Ingeniería de Sistemas y Computación}\\
    {\Large Calidad de Productos de Software}\\
    {\Large 750029C - 01}\\[2.5cm]
    
    {\LARGE \textbf{Sustentación 3}}\\[0.3cm]
    {\Large Pruebas No Funcionales}\\[2cm]
    
    {\large \textbf{Integrantes:}}\\
    \vspace{0.2cm}
    Estiven Andrés Martinez Granados, 2179687\\
    Juan David Loaiza Santiago, 2177570\\
    Jose Otoniel Henao Montes, 2181039\\
    Juan Sebastian Muñoz Rojas, 2177436\\
    Julián David Rendón Cardona, 2177387\\[2cm]
    
    {\large \textbf{Profesora:}}\\
    {\large Beatriz Eugenia Florián Gaviria}\\[2cm]
    
    {\large Santiago de Cali, \today}
    
    \vfill
\end{titlepage}

% ============================================
% TABLA DE CONTENIDO
% ============================================
\tableofcontents
\newpage

% ============================================
% 1. INTRODUCCIÓN
% ============================================
\section{Introducción}

Este informe documenta la validación no funcional de la plataforma Odoo ERP \& CRM v19.0 Community, específicamente en los módulos de \textbf{Ventas} y \textbf{Productos}. Las pruebas se enfocan en dos factores de calidad: \textbf{Eficiencia} y \textbf{Seguridad}.

Para la ejecución de las pruebas, se utilizaron herramientas como \textbf{Apache JMeter} para pruebas de carga, \textbf{Lighthouse} para análisis de rendimiento de interfaz, y pruebas \textbf{manuales en Odoo} para validar controles de acceso y permisos.

El código fuente de las pruebas automatizadas se encuentra disponible en el repositorio de GitHub: \\
\url{https://github.com/estiradikal/odooNF}

\newpage

% ============================================
% 2. DISEÑO DE PRUEBAS NO FUNCIONALES
% ============================================
\section{Diseño de Pruebas No Funcionales}

\subsection{Factores y Criterios de Calidad Seleccionados}

Se seleccionaron los siguientes factores de calidad con sus respectivos criterios:

\begin{table}[H]
\centering
\caption{Factores y criterios de calidad}
\begin{tabular}{|p{4cm}|p{8cm}|}
\hline
\textbf{Factor de Calidad} & \textbf{Criterios de Calidad} \\
\hline
Eficiencia & Tiempo de respuesta, Comportamiento bajo carga, Tiempo de respuesta de interfaz \\
\hline
Seguridad & Control de acceso a productos, Control de permisos de edición \\
\hline
\end{tabular}
\end{table}

\subsection{Métricas e Instrumentos de Inspección e Indagación}

\subsubsection{Métricas seleccionadas}
\begin{itemize}
    \item \textbf{Tiempo promedio de respuesta (ms):} Para medir el tiempo de carga del módulo de Ventas.
    \item \textbf{Tasa de error (\%):} Para medir el comportamiento del sistema bajo carga.
    \item \textbf{Puntuación de Performance (0-100):} Para medir el rendimiento de la interfaz con Lighthouse.
    \item \textbf{¿El sistema bloquea el acceso/edición? (Sí/No):} Para evaluar controles de seguridad.
\end{itemize}

\subsubsection{Instrumentos de Inspección}
\begin{itemize}
    \item \textbf{JMeter:} Herramienta de código abierto para pruebas de carga y rendimiento. Se utilizó para simular usuarios concurrentes y medir tiempos de respuesta y tasa de error.
    \item \textbf{Lighthouse:} Herramienta integrada en Chrome DevTools para analizar el rendimiento, accesibilidad y buenas prácticas de la interfaz web.
    \item \textbf{Pruebas manuales en Odoo:} Para validar controles de acceso y permisos de seguridad.
\end{itemize}

\subsubsection{Instrumentos de Indagación}
No se utilizaron instrumentos de indagación (encuestas, entrevistas) ya que las pruebas no funcionales realizadas son de naturaleza técnica y no requieren intervención de usuarios finales.

\subsection{Matriz de Requerimientos de Pruebas No Funcionales (MRP-NF)}

\begin{longtable}{|p{1.2cm}|p{2cm}|p{2cm}|p{2.5cm}|p{4cm}|p{2.2cm}|p{2.5cm}|p{2.5cm}|p{1.5cm}|p{3cm}|}
\hline
\textbf{ID} & \textbf{Módulo} & \textbf{Factor} & \textbf{Criterio} & \textbf{Escenario} & \textbf{Herramienta} & \textbf{Métrica} & \textbf{Resultado Esperado} & \textbf{Estado} & \textbf{Resultado Obtenido} \\
\hline
\endfirsthead
\hline
\textbf{ID} & \textbf{Módulo} & \textbf{Factor} & \textbf{Criterio} & \textbf{Escenario} & \textbf{Herramienta} & \textbf{Métrica} & \textbf{Resultado Esperado} & \textbf{Estado} & \textbf{Resultado Obtenido} \\
\hline
\endhead
\hline
\multicolumn{10}{r}{Continúa en la siguiente página}
\endfoot
\hline
\endlastfoot
% ===== DATOS DE LA MATRIZ =====
NF-01 & Ventas & Eficiencia & Tiempo de respuesta & Cargar módulo de Ventas con 10, 25 y 50 usuarios concurrentes & JMeter & Tiempo promedio de respuesta (ms) & < 3000 ms & Ok & 287.3 / 190.3 / 239.2 ms (todos < 3000 ms, 0\% error) \\
\hline
NF-02 & Ventas & Eficiencia & Comportamiento bajo carga & Simular 10, 25 y 50 usuarios creando cotizaciones simultáneamente & JMeter & Tasa de error (\%) & 0\% & Ok & 0\% en 10, 25 y 50 usuarios \\
\hline
NF-03 & Productos & Eficiencia & Tiempo de respuesta de interfaz & Analizar el rendimiento de la página principal y módulo de Ventas & Lighthouse & Puntuación de Performance (0-100) & > 90 & NoK & 49/100 (FCP: 5.5s, LCP: 7.0s, TBT: 200ms, Speed Index: 5.5s) \\
\hline
NF-04 & Productos & Seguridad & Control de acceso a productos & Intentar acceder al módulo de Productos con un usuario sin permisos (rol de solo lectura) & Manual (Odoo) & ¿El sistema bloquea el acceso? (Sí/No) & Acceso denegado & NoK & El sistema permitió editar el producto sin los permisos necesarios \\
\hline
NF-05 & Productos & Seguridad & Control de permisos de edición & Intentar editar un producto con un usuario sin permisos de escritura & Manual (Odoo) & ¿El sistema bloquea la edición? (Sí/No) & Acción denegada (campos bloqueados) & NoK & El sistema permitió la edición y guardado del precio \\
\hline
\end{longtable}

\subsection{Selección del Software de Automatización}

El software seleccionado para automatizar el proceso de pruebas es:

\begin{itemize}
    \item \textbf{Apache JMeter:} Para pruebas de carga y rendimiento (NF-01 y NF-02).
    \item \textbf{Lighthouse:} Para análisis de rendimiento de interfaz (NF-03).
    \item \textbf{Odoo (Manual):} Para pruebas de seguridad (NF-04 y NF-05).
\end{itemize}

\newpage

% ============================================
% 3. EJECUCIÓN DE PRUEBAS Y EVIDENCIAS
% ============================================
\section{Ejecución de Pruebas y Evidencias}

\subsection{NF-01: Tiempo de respuesta del módulo Ventas}

Se utilizó JMeter con un Thread Group configurado para 10, 25 y 50 usuarios concurrentes. La URL probada fue \texttt{http://localhost:8069/odoo/sales}.

\begin{table}[H]
\centering
\caption{Resultados NF-01}
\begin{tabular}{|c|c|c|c|c|}
\hline
\textbf{Usuarios} & \textbf{Avg (ms)} & \textbf{Min (ms)} & \textbf{Max (ms)} & \textbf{Error \%} \\
\hline
10 & 287.3 & 115 & 1309 & 0\% \\
25 & 190.3 & 104 & 391 & 0\% \\
50 & 239.2 & 104 & 735 & 0\% \\
\hline
\end{tabular}
\end{table}

\textbf{Evidencia:} Los resultados detallados se encuentran en los archivos anexos \texttt{resultados\_10\_ventas.jtl}, \texttt{resultados\_25\_ventas.jtl} y \texttt{resultados\_50\_ventas.jtl}.

\subsection{NF-02: Comportamiento bajo carga en cotizaciones}

Se utilizó JMeter para simular la creación de cotizaciones a través de la API JSON-RPC de Odoo.

\begin{table}[H]
\centering
\caption{Resultados NF-02}
\begin{tabular}{|c|c|c|}
\hline
\textbf{Usuarios} & \textbf{Avg (ms)} & \textbf{Tasa de error (\%)} \\
\hline
10 & 1058 & 0\% \\
25 & 144 & 0\% \\
50 & 116 & 0\% \\
\hline
\end{tabular}
\end{table}

\textbf{Evidencia:} Los resultados detallados se encuentran en los archivos anexos \texttt{resultados\_api\_10.jtl}, \texttt{resultados\_api\_25.jtl} y \texttt{resultados\_api\_50.jtl}.

\subsection{NF-03: Análisis de rendimiento con Lighthouse}

Se ejecutó Lighthouse en la página principal de Odoo (\texttt{http://localhost:8069}).

\begin{table}[H]
\centering
\caption{Resultados NF-03}
\begin{tabular}{|c|c|}
\hline
\textbf{Métrica} & \textbf{Resultado} \\
\hline
Performance Score & 49/100 \\
FCP & 5.5 s \\
LCP & 7.0 s \\
TBT & 200 ms \\
CLS & 0 \\
Speed Index & 5.5 s \\
\hline
\end{tabular}
\end{table}

\textbf{Evidencia:} El reporte completo de Lighthouse se encuentra en el anexo \texttt{reporte\_lighthouse.html}.

\subsection{NF-04: Control de acceso a productos}

Se creó un usuario con permisos limitados y se intentó acceder al módulo de Productos.

\begin{table}[H]
\centering
\caption{Resultados NF-04}
\begin{tabular}{|c|c|}
\hline
\textbf{Métrica} & \textbf{Resultado} \\
\hline
¿El sistema bloquea el acceso? & No \\
\hline
Resultado & El sistema permitió editar el producto sin los permisos necesarios \\
\hline
\end{tabular}
\end{table}

\textbf{Evidencia:} La captura de pantalla se encuentra en el anexo \texttt{NF-04-acceso-denegado.png}.

\subsection{NF-05: Control de permisos de edición}

Se creó un usuario con permisos de solo lectura y se intentó editar un producto.

\begin{table}[H]
\centering
\caption{Resultados NF-05}
\begin{tabular}{|c|c|}
\hline
\textbf{Métrica} & \textbf{Resultado} \\
\hline
¿El sistema bloquea la edición? & No \\
\hline
Resultado & El sistema permitió la edición y guardado del precio \\
\hline
\end{tabular}
\end{table}

\textbf{Evidencia:} La captura de pantalla se encuentra en el anexo \texttt{NF-05-campos-bloqueados.png}.

\newpage

% ============================================
% 4. ANÁLISIS DE RESULTADOS
% ============================================
\section{Análisis de Resultados}

\subsection{Evaluación e interpretación de las métricas}

\subsubsection{NF-01: Tiempo de respuesta}

Los tiempos promedio de respuesta para cargar el módulo de Ventas fueron de \textbf{287.3 ms} (10 usuarios), \textbf{190.3 ms} (25 usuarios) y \textbf{239.2 ms} (50 usuarios). Todos los valores están muy por debajo del umbral de 3000 ms, lo que indica un rendimiento excelente. La tasa de error fue del 0\% en todos los casos.

\subsubsection{NF-02: Comportamiento bajo carga}

La tasa de error fue del 0\% en todos los escenarios. El tiempo promedio de creación de cotizaciones vía API fue de \textbf{1058 ms} (10 usuarios), \textbf{144 ms} (25 usuarios) y \textbf{116 ms} (50 usuarios). El sistema es estable y eficiente bajo carga moderada. El rendimiento mejora al aumentar la carga, lo que sugiere que el sistema aprovecha mejor los recursos con más conexiones.

\subsubsection{NF-03: Rendimiento de interfaz}

Lighthouse reportó una puntuación de \textbf{49/100} en Performance, lo que indica que la interfaz de Odoo tiene problemas de optimización. Los tiempos de carga (LCP: 7.0s, FCP: 5.5s) son excesivamente altos y afectan negativamente la experiencia del usuario.

\subsubsection{NF-04: Control de acceso}

El sistema permitió a un usuario sin permisos editar un producto. Esto representa una vulnerabilidad de seguridad que debe ser corregida.

\subsubsection{NF-05: Control de permisos de edición}

El sistema permitió la edición y guardado del precio de un producto a un usuario sin permisos de escritura. Esta es una vulnerabilidad de seguridad grave que afecta la integridad de los datos.

\newpage

% ============================================
% 5. REPORTE DE NO CONFORMIDADES
% ============================================
\section{Reporte de No Conformidades}

\subsection{Resumen de ejecución}

\begin{table}[H]
\centering
\caption{Resumen de ejecución}
\begin{tabular}{|l|c|}
\hline
\textbf{Métrica de Ejecución} & \textbf{Cantidad} \\
\hline
Total de Requerimientos de Prueba Diseñados & 5 \\
Total de Requerimientos Ejecutados & 5 \\
Requerimientos Exitosos (Ok) & 2 \\
Requerimientos con Fallos (NoK) & 3 \\
No Conformidades Detectadas & 3 \\
\hline
\end{tabular}
\end{table}

\subsection{Reporte de No Conformidades}

\begin{table}[H]
\centering
\caption{No Conformidades Detectadas}
\begin{tabular}{|p{1.5cm}|p{2.5cm}|p{3.5cm}|p{3.5cm}|p{2.5cm}|p{2cm}|}
\hline
\textbf{ID} & \textbf{Dónde} & \textbf{Cuándo} & \textbf{Hallazgo} & \textbf{Se Espera} & \textbf{Severidad} \\
\hline
NC-03 & Módulo de Productos & Al ejecutar Lighthouse en la página principal & Performance Score bajo (49/100), FCP 5.5s, LCP 7.0s & Performance Score > 90, tiempos de carga óptimos & Eficiencia (Rendimiento) \\
\hline
NC-04 & Módulo de Productos & Al intentar acceder con usuario sin permisos & El sistema permitió editar el producto sin permisos & El sistema debe bloquear el acceso a la edición & Seguridad \\
\hline
NC-05 & Módulo de Productos & Al intentar editar con usuario sin permisos de escritura & El sistema permitió la edición y guardado del precio & El sistema debe bloquear la edición de campos para usuarios sin permisos & Seguridad \\
\hline
\end{tabular}
\end{table}

\newpage

% ============================================
% 6. LECCIONES APRENDIDAS
% ============================================
\section{Lecciones Aprendidas}

\subsection{Aspectos Positivos (a repetir)}

\begin{itemize}
    \item La automatización con \textbf{JMeter} permitió simular carga de manera eficiente y repetible, facilitando la obtención de métricas precisas.
    \item \textbf{Lighthouse} proporcionó métricas claras y útiles para identificar problemas de rendimiento en la interfaz.
    \item El entorno \textbf{Docker} garantizó la reproducibilidad de las pruebas y un ambiente controlado.
    \item La estructura de la matriz de requerimientos (MRP-NF) permitió una trazabilidad clara de cada caso de prueba.
\end{itemize}

\subsection{Aspectos a Mejorar}

\begin{itemize}
    \item Configurar JMeter para manejar autenticación y sesiones de forma más robusta en pruebas de API.
    \item Incluir más pruebas de seguridad con roles intermedios y permisos granulares para identificar todas las vulnerabilidades.
    \item Realizar pruebas con cargas superiores a 100 usuarios para identificar el límite de escalabilidad del sistema.
    \item Optimizar el frontend de Odoo para reducir los tiempos de carga y mejorar la experiencia del usuario.
    \item Revisar la configuración de permisos y roles en Odoo para garantizar que los usuarios con permisos limitados no puedan modificar información sensible.
\end{itemize}

\newpage

% ============================================
% 7. ANEXOS
% ============================================

A continuación, se listan los archivos complementarios que soportan las evidencias presentadas en este informe.

\subsection*{Repositorio de GitHub}

El código fuente de las pruebas automatizadas y los planes de JMeter se encuentran disponibles en el repositorio de GitHub: \\
\url{https://github.com/estiradikal/odooNF}

\subsection*{Matriz de Requerimientos de Pruebas en Excel}

La matriz de requerimientos de pruebas no funcionales (MRP-NF) con todos los casos diseñados y resultados obtenidos se encuentra disponible en el siguiente enlace de Google Sheets: \\
\url{https://docs.google.com/spreadsheets/d/1THcPloRmVX_CbkE1HrzkfvkAeIVdkhuUGCu59apmvMU/edit?gid=1160611668#gid=1160611668}

\end{document}